\documentclass[11pt]{article}
\usepackage{fontspec}
\setmainfont{Latin Modern Roman}
\usepackage{amsmath,amssymb,mathtools,unicode-math}
\setmathfont{Latin Modern Math}
\usepackage{xurl,hyperref}
\hypersetup{hidelinks,pdftitle={Problem statement}}
\newcommand{\sref}[2]{\hyperlink{source:#1:#2}{[S#2]}}
\newcommand{\cataloguescope}[1]{\paragraph{Mathematical question.}#1}
\begin{document}
\section{Open Problems in Integrable Systems — Polynomial in momenta integrals in higher dimensions.}\label{statement}

\subsection{Definitions and mathematical statement}

\paragraph{Shared notation.}$\mathbb N=\{1,2,\ldots\}$, and $\mathbb Z,\mathbb Q,\mathbb R,\mathbb C$ denote the integers, rationals, reals and complex numbers. Any other conventions are specified locally.


Let $(M,g)$ be a smooth Riemannian manifold, with curvature tensor $R$ and Levi-Civita connection $\nabla$. A scalar invariant of order at most one is a coordinate-independent algebraic expression formed from $R$, $\nabla R$ and the metric, for example a complete contraction of their tensor products. The source calls $(M,g)$ $1$-homogeneous when all such scalar invariants are constant on $M$. It is locally homogeneous if the local isometry pseudogroup acts transitively, equivalently any two sufficiently nearby points are related by a local isometry.
\paragraph{Statement recorded in the supplied source.}
Is every $1$-homogeneous Riemannian manifold locally homogeneous? \sref{statement}{2}
\subsection{Short English statement}
Does constancy of all first-order scalar curvature invariants imply local homogeneity for positive-definite metrics?
\subsection{Sources}
\begin{description}
\item[\hypertarget{source:statement:1}{S1}] ULAM.AI and the UnsolvedMath Contributors, UnsolvedMath: A Curated Collection of Open Mathematics Problems, catalogue entry. CC BY 4.0. \url{https://huggingface.co/datasets/ulamai/UnsolvedMath}.
\item[\hypertarget{source:statement:2}{S2}] A. Bolsinov, V. S. Matveev, E. Miranda and S. Tabachnikov, Open Problems, Questions, and Challenges in Finite-Dimensional Integrable Systems (2018), arXiv:1804.03737, Question 3.8. \url{https://arxiv.org/abs/1804.03737}.
\end{description}
\label{end:statement}
\end{document}
